% ============================================================
%  计算机软件著作权登记申请表（正文）
% ============================================================
% 用法：直接修改下面各字段的第二个参数即可。勾选项用
%   \tjcrchk{...}    已勾选
%   \tjcrunchk{...}  未勾选
% 需要打印后手写的空白处用 \tjcrblank{宽度}。
% 字数限制等提示文字用 \tjcrred{...} 标红。
% “软件的主要功能”字数较多，已按句号/分号拆成若干续行
% （\tjcrcont），以便 longtable 在分页处自然断开；增删内容时
% 保持每条续行不超过本页高度即可。
% ============================================================

\vspace*{0.2cm}
{\centering\linespread{1}\fontsize{12bp}{18bp}\selectfont\heiti\bfseries
  计算机软件著作权登记申请表\par}
\vspace{0.5em}

\begin{tjcrform}

% ------------------------------------------------------------
% 权利取得方式
% ------------------------------------------------------------
\tjcrrow{权利取得方式}{%
  \linespread{1.5}\selectfont
  \tjcrchk{原始取得}\newline
  \tjcrunchk{继受取得}\quad(\tjcrunchk{受让}\quad\tjcrunchk{承受}\quad\tjcrunchk{继承})\newline
  \hspace*{2em}\tjcrunchk{软件是否已登记}\quad（原登记号：\tjcrblank{3cm}）\newline
  \hspace*{2em}\tjcrunchk{原登记是否做过变更或补充}\newline
  \hspace*{4em}（变更或补充证明编号：\tjcrblank{2.8cm}）}

% ------------------------------------------------------------
% 软件基本信息
% ------------------------------------------------------------
\tjcrsection{软件基本信息}

\tjcrrow{软件名称}{码极客具身导航评测与诊断软件}
\tjcrrow{软件简称}{码极客 NavEval Studio}
\tjcrrow{版本号}{V1.0}

\tjcrrow{权利范围}{%
  \linespread{1.5}\selectfont
  \tjcrchk{全部}\newline
  \tjcrunchk{部分}\newline
  \hspace*{2em}\tjcrunchk{发表权}\quad\tjcrunchk{署名权}\quad\tjcrunchk{修改权}\quad
  \tjcrunchk{复制权}\quad\tjcrunchk{发行权}\newline
  \hspace*{2em}\tjcrunchk{出租权}\quad\tjcrunchk{信息网络传播权}\quad
  \tjcrunchk{翻译权}\quad\tjcrunchk{应当由著作权人享有的其他权利}}

% ------------------------------------------------------------
% 软件开发信息
% ------------------------------------------------------------
\tjcrsection{软件开发信息}

\tjcrrow{软件分类}{%
  \tjcrchk{应用软件}\quad\tjcrunchk{嵌入式软件}\quad
  \tjcrunchk{中间件}\quad\tjcrunchk{操作系统}}

\tjcrrow{软件作品说明}{%
  \linespread{1.5}\selectfont
  \tjcrchk{原创}\qquad\tjcrunchk{修改（含翻译软件、合成软件）}\newline
  \hspace*{4em}修改合成或翻译说明（50字内）：\newline
  \tjcrunchk{该软件已登记}\quad\tjcrblank{3.2cm}\quad（上传原证书扫描件）\newline
  是修改、翻译或合成别人软件且需要原权利人授权的（上传授权书）}

\tjcrrow{开发方式}{%
  \linespread{1.5}\selectfont
  \tjcrchk{独立开发}\quad\tjcrunchk{合作开发}\quad\tjcrunchk{委托开发}\quad
  \tjunchecked\hspace{0.25em}下达任务开发（除独立开发，其他开发方式需要提供合同或协议书）}

\tjcrrow{开发完成日期}{2026.11.09}
\tjcrrow{发布状态}{\tjcrunchk{已发表}\qquad\tjcrchk{未发表}}

% “已发表信息”两小行放在同一单元格内，行间不画横线，标签自然居中。
\tjcrrow{已发表信息}{%
  \linespread{1.5}\selectfont
  首次发表日期：\tjcrblank{3cm}\newline
  首次发表地点：国家：中国\quad 城市：上海}

% ------------------------------------------------------------
% 著作权人（子表：表头 + 两行空白填写行）
% ------------------------------------------------------------
\tjcrrowraw{\tjcrstrut 著作权人\newline（业务部门\newline 只需填写\newline 姓名或名称）}{%
  \begin{tjcrcopyrighters}
    姓名或名称 & 类别 & 证件类型 & 证件号码 & 国籍 & 省份/\newline 城市 & 园区 \\ \hline
    & & & & & & \\ \hline
    & & & & & & \\ \hline
  \end{tjcrcopyrighters}%
  % 略增本行高度，使“软件功能与特点”分节行另起一页，避免孤行。
  \vspace*{0.5\baselineskip}}

% ------------------------------------------------------------
% 软件功能与特点
% ------------------------------------------------------------
\tjcrsection{软件功能与特点}

\tjcrrow{开发的硬件环境（\tjcrred{50字符内}）}%
  {x86-64处理器、8GB以上内存的台式机或笔记本}
\tjcrrow{运行的硬件环境（\tjcrred{50字符内}）}%
  {x86\_64架PC或笔记本电脑，Intel Core i5及以上处理器，8GB及以上内存}
\tjcrrow{开发该软件的操作系统（\tjcrred{50字符内}）}{Ubuntu 20.04 LTS}
\tjcrrow{软件开发环境/开发工具（\tjcrred{50字符内}）}%
  {VS Code、Python、Vue 3、Chrome开发者工具}
\tjcrrow{该软件的运行平台/操作系统（\tjcrred{50字符内}）}%
  {Ubuntu、Windows，配合Chrome或Edge浏览器}
\tjcrrow{软件运行支撑环境/支持软件（\tjcrred{50字符内}）}%
  {Python3.9以上、FastAPI、SQLite，可选habitat\_sim}

\tjcrrow{编程语言}{%
  \linespread{1.5}\selectfont
  \tjcrunchk{Assembly language}\quad\tjcrunchk{C}\quad\tjcrunchk{C\#}\quad
  \tjcrunchk{C++}\quad\tjcrunchk{Delphi/Object Pascal}\quad\tjcrunchk{Go}\quad
  \tjcrchk{HTML}\newline
  \tjcrunchk{Java}\quad\tjcrchk{JavaScript}\quad\tjcrunchk{MATLAB}\quad
  \tjcrunchk{Objective-C}\quad\tjcrunchk{PHP}\quad\tjcrunchk{PL/SQL}\quad
  \tjcrunchk{Perl}\quad\tjcrchk{Python}\quad\tjcrunchk{R}\quad
  \tjcrunchk{Ruby}\quad\tjcrchk{SQL}\quad\tjcrunchk{Swift}\newline
  \tjcrunchk{Visual Basic}\quad\tjcrunchk{Visual Basic .Net}\newline
  如有需要，请输入其他编程语言：}

\tjcrrow{源程序量}{25600行}
\tjcrrow{开发目的（\tjcrred{50字符内}）}%
  {为连续环境视觉语言导航模型提供可视化评测、诊断与实时仿真工具。}
\tjcrrow{面向领域/行业（\tjcrred{50字符内}）}%
  {人工智能、具身智能、机器人导航与高校科研。}

% ------------------------------------------------------------
% 软件的主要功能：按句号/分号拆成多条续行，便于跨页
% ------------------------------------------------------------
\tjcrrow{软件的主要功能\newline（\tjcrred{500-1300字内}）}{%
  码极客具身导航评测与诊断软件面向连续环境视觉语言导航模型，
  提供从数据准备、任务调度、过程监控、结果分析到轨迹回放与实时仿真的一体化工程环境。}
\tjcrcont{%
  数据与配置方面，软件支持数据集登记与扫描、Episode 检索、多源观测时间轴、
  数据质量校验、全集与冒烟集切分生成，以及严格时间前缀的因果切片导出。}
\tjcrcont{%
  评测调度方面，软件提供配置模板与模型档案管理、模型可用性探测、任务队列与并发控制、
  GPU 互斥、事件流实时监控、暂停恢复取消与失败重跑，并周期性采集显存、时延、上下文长度等资源数据。}
\tjcrcont{%
  推理与回放方面，软件以适配器协议对接外部推理后端，调度真实模型时逐帧采集位姿、朝向、动作、
  单步时延、记忆窗口与裁剪前后 token 数；回放工作台四区联动，同屏显示第一视角画面、
  俯视轨迹与参考路径、动作时间轴和记忆状态曲线，支持拖动定位、单步、倍速与循环播放。}
\tjcrcont{%
  分析诊断方面，软件聚合成功率、加权成功率、导航误差、oracle 成功率、归一化动态时间规整与轨迹长度，
  提供多任务对比与基线差值、模态与记忆策略消融矩阵、配对 bootstrap 显著性检验、失败归因、
  长时性能曲线，并可导出 Markdown 与 CSV 报告。}
\tjcrcont{%
  交互式仿真方面，软件可从已登记 Episode 启动单局仿真，第一视角以视频流实时推送，
  俯视图同时绘制实际轨迹、人类示范轨迹与最短路径，支持手动、模型驾驶与 oracle 自动驾驶三种方式，
  并提供第三视角画中画、机器狗模型与会话录制导出。}
\tjcrlastcont{%
  系统与资产方面，软件提供用户角色与操作审计、系统设置与 GPU 状态监测，
  可自动导入备份库中的数据集、场景、模型、历史评测结果与资料文件，
  并在无 GPU 环境下以离线演示模式跑通全部功能链路。}

% ------------------------------------------------------------
% 软件的技术特点
% ------------------------------------------------------------
\tjcrrow{软件的技术特点\newline（\tjcrred{最多3个，标签可选}）}{%
  \linespread{1.5}\selectfont
  \tjcrunchk{APP}\quad\tjcrunchk{游戏软件}\quad\tjcrunchk{教育软件}\quad
  \tjcrunchk{金融软件}\quad\tjcrunchk{医疗软件}\quad\tjcrunchk{地理信息软件}\newline
  \tjcrunchk{云计算软件}\quad\tjcrunchk{信息安全软件}\quad
  \tjcrunchk{大数据软件}\quad\tjcrchk{人工智能软件}\quad
  \tjcrunchk{VR软件}\quad\tjcrunchk{5G软件}\quad\tjcrunchk{小程序}\quad
  \tjcrunchk{物联网软件}\quad\tjcrunchk{智慧城市软件}\par
  \vspace{0.4em}
  技术特点（100字内）：（必填）适配器协议解耦模型，逐帧元数据打通评测与回放，
  配置快照保证复现，输入白名单保证因果性，多线程解耦实现模型在环仿真。}

\end{tjcrform}
